% SPDX-License-Identifier: CC-BY-SA-4.0
% Author: Matthieu Perrin
% Part: <Nom de la partie>
% Section: <Nom de la section>
% Sub-section: <Nom de la sous-section>  % (facultatif, laisser vide si non utilisé)
% Frame: <Titre de la slide>

\begingroup

\begin{frame}{Encodage effectif et langage universel}

  \vspace{-2mm}
  \begin{block}{Définition -- Encodage effectif des MTD}
    Un \structure{encodage effectif} des MTD est un formalisme 
    \alert{$\langle \mathcal{L}, \llbracket \cdot \rrbracket \rangle$} tel que :
    \begin{description}[Turing-complet :]
    \item[Turing-complet :] $\alert{\operatorname{Im}(\llbracket \cdot \rrbracket) = \textsc{re}}$
      \begin{itemize}\small
      \item Tout problème semi-décidable est \structure{encodable} dans $\mathcal{L}$
      \item Les MTDs sont encodées à \structure{équivalence sémantique} près
      \item $\llbracket \cdot \rrbracket$ n'est pas forcément bijective
      \end{itemize}
    \item[Effectif :] $\alert{\left\{ \langle P, u \rangle \mid P \in \mathcal{L} \land u \in \llbracket P \rrbracket \right\} \in \textsc{re}}$
      \begin{itemize}\small
      \item On peut décoder les MTD et les simuler
      \item Il s'agit du langage universel :
      \end{itemize}
    \end{description}
        $$\structure{\textsc{Universal}_{\langle \mathcal{L}, \llbracket \cdot \rrbracket \rangle} \eqdef \left\{ \langle P, u \rangle \mid P \in \mathcal{L} \land u \in \llbracket P \rrbracket \right\}}$$

  \end{block}

  \pause

  \vspace{-2mm}
  \begin{alertblock}{Remarques -- Langage de programmation}
    Un \alert{langage de programmation} est un encodage effectif des MTD
    \begin{itemize}
    \item On appellera les mots de $\mathcal{L}$ des \alert{programmes}
    \item L'effectivité signifie qu'il existe un \alert{interpréteur} pour le langage
    \end{itemize}

    \begin{center}
      \alert{$\textsc{Universal}_{\langle \mathcal{L}, \llbracket \cdot \rrbracket \rangle}$ est-il décidable ?}
    \end{center}
  \end{alertblock}

  
\end{frame}
